% 四、本课题的特色及创新之处。
% 内容迁自 NSFC_2026_MineChangeTrust；保留原稿事实、暂定值和待补项。
\subsection{项目特色}
本项目同时考察矿山变化检测的误报和细小变化漏检，比较一致性约束作用于不同区域时的效果，分析变化尺度和外观差异对结果的影响。研究采用公开数据，将不同用途的数据按矿区隔离，并完整报告失效矿区和反例。对总体指标的改善，还需进一步查明其成立条件能否重复。

可靠性评价在检测器选定且参数固定后开展，分别记录概率误差、复核面积和全图剩余错误。通过保持检测结果一致，区分检测改进与复核选择各自带来的收益，并查明哪些高置信度漏检仍未被发现。

\subsection{创新点}
\textbf{（1）区分约束区域、强度与判决工作点对检出权衡的影响。}
预期贡献是查明背景限定约束在何种条件下缓解误报与小变化召回的冲突。同骨干区域对照、强度匹配和阈值敏感性分析分别检验空间选择、正则程度及分数尺度的解释，固定模型的污染掩码提供梯度证据。相对全图约束改善、相对无约束产生增量、在DonaNet上可迁移分别报告；仅减轻全图约束损伤时不称为检测方法突破。机制认识须由排除替代解释的对照支持。

\textbf{（2）检验可靠性改善能否帮助识别约束后残留的错误。}
在冻结检测器与相同面积预算下，分析概率误差降低是否伴随误报、漏检排序改善，重点检查高置信度漏检。复用无约束、全图和背景组的预测，追踪复核仍无法发现的错误，解释检测权衡与复核局限的关系。温度缩放、熵及多视图分歧均为已有工具；科学增量在于建立概率校准、错误可识别性与全图剩余风险之间的实证联系，离线结果限于理想纠错上限。
